\chapter*{Abstract}
\addcontentsline{toc}{chapter}{Abstract}

Our group has built a C++ toolkit called \textbf{di3p} that enhances images by combining spatial
segmentation with frequency-domain filtering. The core idea is straightforward: instead
of applying the same filter to an entire image, we first use K-Means clustering to split
the image into colour-coherent regions, and then treat each region independently in the
frequency domain using the Discrete Fourier Transform (DFT) and the Discrete Cosine
Transform (DCT). This lets us apply targeted low-pass or high-pass filters to specific
parts of the scene without touching the rest. The system was tested on our own photographs, using Shannon entropy
reduction and PSNR as metrics to confirm that noise was being suppressed
without degrading structure. 

\vspace{0.5cm}
\noindent The entire pipeline was written from scratch in
C++17 with no dependency on OpenCV or similar libraries.

\vspace{0.5cm}
\noindent\textbf{Keywords:} image segmentation, K-Means clustering, Discrete Fourier
Transform, Discrete Cosine Transform, frequency-domain filtering, image enhancement,
PSNR, C++, digital image processing.
